% Figura extraída de "Tema 3.tex" (línea 60 del original)
\documentclass[12pt,border=1pt]{standalone}
\usepackage[spanish,es-tabla]{babel}
\usepackage[utf8]{inputenc}
\usepackage{amsmath}
\usepackage{xcolor}
\usepackage{tikz}
\usepackage{pgfplots}
\usepackage{circuitikz}
\usetikzlibrary{babel, circuits.ee}

\definecolor{applegreen}{rgb}{0.55, 0.71, 0.0}
\definecolor{airforceblue}{rgb}{0.36, 0.54, 0.66}
\definecolor{amber}{rgb}{1.0, 0.49, 0.0}

\newcommand{\diff}[3][1]{%
	\IfBlankTF{#1}{\frac{\mathrm{d}{#2}}{\mathrm{d}{#3}}}%
	{\frac{\mathrm{d}^{#1} {#2}}{\mathrm{d}{#3}^{#1}}}%
}%
\newcommand{\pdiff}[3][1]{%
	\IfBlankTF{#1}{\frac{\partial {#2}}{\partial {#3}}}%
	{\frac{\partial^{#1} {#2}}{\partial {#3}^{#1}}}%
}%

\begin{document}
\begin{tikzpicture}
  % Cuadrado exterior
  \draw (0,0) -- (4,0) -- (4,4) -- (-4,4) -- (-4,0) -- cycle;
  \draw (0,0) -- (-4,0) -- (-4,4) -- (0,4) -- cycle;
  \draw (-6,0) rectangle (-4,4);
  \draw (6,0) rectangle (4,4);
  \draw [dashed] (-4,4)--(-4,-3);
  \draw [dashed] (4,4)--(4,-3);
  
  \draw[latex -latex] (-4,-0.5) -- (4,-0.5);
  \node at (0,-0.8) {Zona de Carga Espacial};
  \draw [-latex, ultra thick] (2,-1.4)--(-2,-1.4);
  \node at (0,-1.7) {\textbf{E}};

  \node at (-5,2) {P};
  \node at (5,2) {N};

  \draw[-latex]  (-4.5,-2.2)--(-4,-2.2);
  \node at (-5.5,-2) {Fuerzas};
  \node at (-5.5,-2.4) {de difusión};
  \node at (-5.5,-2.8) {en $h^+$};
  \draw[-latex]  (4.5,-2.2)--(4,-2.2);
  \node at (5.5,-2) {Fuerzas};
  \node at (5.5,-2.4) {de difusión};
  \node at (5.5,-2.8) {en $e^-$};

  \draw[-latex]  (-3.5,-2.2)--(-4,-2.2);
  \node at (-2.5,-2) {Fuerzas};
  \node at (-2.5,-2.4) {eléctricas};
  \node at (-2.5,-2.8) {en $h^+$};
  \draw[-latex]  (3.5,-2.2)--(4,-2.2);
  \node at (2.5,-2) {Fuerzas};
  \node at (2.5,-2.4) {eléctricas};
  \node at (2.5,-2.8) {en $e^-$};
  % Cuadrado interior 4x4 con signos "-"
  \foreach \x in {-3,-2,-1,0} {
\foreach \y in {0,1,2,3} {
  \node at (\x-0.5, \y+0.5) {$-$};
}
  }
  % Cuadrado interior 4x4 con signos "+"
  \foreach \x in {1,2,3,4} {
\foreach \y in {1,2,3,4} {
  \node at (\x-0.5, \y-0.5) {+};
}
  }
  \node at (-2,2) {$q(N_\mathrm{A})<0$};
  \node at (2,2) {$q(N_\mathrm{D})>0$};
\end{tikzpicture}
\end{document}
